% Generated draft: requires mathematical and source review.
\documentclass[11pt]{article}
\usepackage[margin=25mm]{geometry}
\usepackage{fontspec}
\setmainfont{Latin Modern Roman}
\usepackage{amsmath,amssymb,mathtools,unicode-math}
\setmathfont{Latin Modern Math}
\usepackage{enumitem,needspace,xurl,hyperref}
\setlength{\parindent}{0pt}
\setlength{\parskip}{5pt}
\setlength{\emergencystretch}{3em}
\newcommand{\N}{\mathbb{N}}
\newcommand{\Z}{\mathbb{Z}}
\newcommand{\Q}{\mathbb{Q}}
\newcommand{\R}{\mathbb{R}}
\newcommand{\C}{\mathbb{C}}
\newcommand{\F}{\mathbb{F}}
\newcommand{\A}{\mathbb{A}}
\newcommand{\PP}{\mathbb{P}}
\newcommand{\E}{\mathbb{E}}
\newcommand{\eps}{\varepsilon}
\newcommand{\dd}{\,\mathrm{d}}
\begin{document}
\section*{UnsolvedMath ID: 30002364}
\textbf{Source problem code:} catalogue entry
\subsection{Definitions and mathematical statement}
Let \(A/\overline{\Q}\) be an abelian variety of dimension \(g\) of CM type, meaning that \(\operatorname{End}(A)\otimes_{\Z}\Q\) contains a commutative semisimple \(\Q\)-subalgebra of dimension \(2g\). Fix an embedding \(\overline{\Q}\hookrightarrow\C\). For each integer \(n\geq 1\), put
\[
\operatorname{Hdg}^{r}(A^n)=H^{2r}(A^n(\C),\Q)\cap H^{r,r}(A^n(\C)),\qquad 0\leq r\leq ng,
\]
where the intersection is taken inside complex cohomology. Let \(L_H^\bullet(A^n)\) be the graded \(\Q\)-subalgebra of even rational cohomology generated by divisor classes. Its elements are called Lefschetz classes. Define
\[
\mathcal H(A):\quad \operatorname{Hdg}^{r}(A^n)=L_H^{r}(A^n)\quad\text{for every }n\geq1\text{ and }0\leq r\leq ng,
\]
where the superscript \(r\) on \(L_H\) denotes codimension, hence cohomological degree \(2r\).

To interpret reduction as an abelian variety, choose a number field \(K\) over which \(A\) is defined and a finite place \(v\) of good reduction, with residue field \(k=\F_q\) of characteristic \(p\). Let \(A_0/k\) be the special fibre. The source does not explicitly specify good reduction; this is the conventional interpretation used here, including good reduction obtained after a finite extension. Bad reduction is not assigned an interpretation in this draft.

For a prime \(\ell\ne p\), let
\[
T_\ell^{r}(A_0^n)=\bigcup_{m\geq1}H_{\mathrm{\acute et}}^{2r}\bigl((A_0^n)_{\overline{k}},\Q_\ell(r)\bigr)^{\operatorname{Gal}(\overline{k}/\F_{q^m})}.
\]
Thus Tate classes here are classes fixed after some finite extension of the residue field. Let \(L_\ell^\bullet(A_0^n)\) be the graded \(\Q_\ell\)-subalgebra of Tate-twisted even \(\ell\)-adic cohomology generated by the cycle classes of divisors on \((A_0^n)_{\overline{k}}\). Define
\[
\mathcal T_\ell(A_0):\quad T_\ell^{r}(A_0^n)=L_\ell^{r}(A_0^n)\quad\text{for every }n\geq1\text{ and }0\leq r\leq ng.
\]
The finite-extension convention for Tate classes and the use of geometric divisor classes are explicit conventions of this draft; the supplied statement does not specify them.

\textbf{Forward question.} Does \(\mathcal H(A)\) imply \(\mathcal T_\ell(A_0)\) for every reduction \(A_0\) of \(A\), interpreted as above, and every prime \(\ell\ne\operatorname{char}(k)\)? The quantifier over every reduction is explicit in the supplied original statement. The quantifier over \(\ell\) makes explicit the usual \(\ell\)-adic reading, but is not stated in the supplied source record.

\textbf{Converse question.} Is there a class of CM-type abelian varieties for which
\[
\mathcal T_\ell(A_0)\ \Longrightarrow\ \mathcal H(A)
\]
holds? The source does not specify the intended class, admissible reduction hypotheses, or whether the antecedent concerns one reduction, every reduction, one auxiliary prime, or every auxiliary prime. These quantifiers remain unresolved; the displayed implication records the requested direction without choosing them.

The supplied dated assessment reports an equivalence proved by Sugiyama for simple CM abelian varieties with an abelian CM field under explicit reduction hypotheses. Those hypotheses are not supplied, so this draft does not state that result as a precise theorem or extend it to arbitrary reductions.
\subsection{Short English statement}
For a CM abelian variety, ask whether generation of all Hodge classes on every power by divisor classes forces the analogous generation of Tate classes on every power of every reduction. Also ask for a class of CM abelian varieties where a converse holds. The reduction conventions and the quantifiers in the converse require clarification. The supplied partially solved label and special-case literature assessment remain dated curation requiring review.
\subsection{Sources}
\begin{itemize}
\item Local UnsolvedMath registry, permanent ID \(30002364\), title \emph{Lefschetz-Class Properties Under CM Reduction}, source problem code \texttt{catalogue entry}; supplied complete source record. Repository registry locations are \texttt{data/unsolvedmath/problems.json} and \texttt{catalog/problems.json}. The record belongs to the Oberwolfach Reports Open Problems collection and gives proposed year 2013; no individual proposer is supplied.
\item Supplied report citation: \emph{Algebraic K-theory and Motivic Cohomology (2014)}, Oberwolfach Reports, pp.\ 1861--1913, DOI 10.4171/owr/2013/32. \url{https://doi.org/10.4171/owr/2013/32}. The citation's 2014 label and the record's 2013 year are retained without reconciliation.
\item Supplied literature source, ``Oberwolfach report on CM reductions'': \url{https://ems.press/content/serial-article-files/46461?nt=1}.
\item Supplied literature source, ``Oberwolfach report on Lefschetz classes'': \url{https://publications.mfo.de/bitstream/handle/mfo/3363/OWR_2013_32.pdf?isAllowed=y&sequence=1}.
\item Supplied literature source, ``Sugiyama, Lefschetz classes'': \url{https://doi.org/10.1016/j.jnt.2012.09.023}.
\item Literature assessment dated exactly \texttt{2026-08-21}: supplied dated curation, not independently verified. It labels the problem partially solved, reports that Sugiyama proves equivalence for simple CM abelian varieties with an abelian CM field under explicit reduction hypotheses, and reports that the questions for arbitrary primes and CM varieties remain open. The record expressly requests independent verification of time-sensitive status and cited claims.
\item The supplied statement status is \texttt{corrected\_verified}. Its verification note reports checking against the official Oberwolfach report and restoration of the source's hypotheses and notation. This is a supplied verification claim, not an independent verification performed for this draft.
\end{itemize}
\end{document}
